% © 2026 Ing. David Jaroš — CC BY-NC-ND 4.0
% UBT-AI-PROVENANCE-BEGIN
% schema: ubt-ai-provenance/v1
% tier: B_machine_verified
% ai_assistance: disclosed
% human_review: machine-verification
% editorial_responsibility: Ing. David Jaroš
% policy: ../../AI_PROVENANCE.md
% notice: Machine-verified against named sources or verifiers; individual attestation is not claimed.
% UBT-AI-PROVENANCE-END

\documentclass[12pt,a4paper]{article}
\usepackage[a4paper,margin=2.5cm]{geometry}
\usepackage{amsmath,amssymb,amsthm,mathtools}
\usepackage[most]{tcolorbox}
\usepackage[hidelinks]{hyperref}
\usepackage{ubtprovenance}
\UBTTier{B}
\newtheorem{definition}{Definition}[section]
\newtheorem{theorem}[definition]{Theorem}
\newtheorem{corollary}[definition]{Corollary}
\theoremstyle{remark}\newtheorem{remark}[definition]{Remark}
\title{GAP-10D: Conditional Low-Energy Uniqueness of Einstein Dynamics}
\author{Ing. David Jaro\v{s}}
\date{16 July 2026}
\begin{document}
\maketitle
\UBTProvenanceNotice

\begin{abstract}
The covariant-tetrad architecture does not by itself derive the Einstein
equations.  This note closes two conditional subproblems.  First, variation of
the minimal Hilbert--Palatini action with respect to the tetrad and independent
connection gives the Einstein equation and the Cartan torsion equation.
Second, Lovelock uniqueness implies that in four dimensions any local natural
symmetric divergence-free metric equation of second differential order has
the form $aG_{\mu\nu}+bg_{\mu\nu}$.  Thus, once UBT derives the standard
low-energy assumptions, the classical metric equation is fixed to Einstein
with a cosmological constant.  The remaining fundamental gap is the derivation
of those assumptions, coefficients and matter coupling from the canonical
single-$\Theta$ action, not ambiguity among many second-order metric equations.
\end{abstract}

\section{Minimal first-order action}
Use the action
\begin{equation}
 S[e,\omega,\Theta]
 =\frac{1}{4\kappa}\int\epsilon_{abcd}
 e^a\wedge e^b\wedge R^{cd}(\omega)
 -\frac{\Lambda}{24\kappa}\int\epsilon_{abcd}
 e^a\wedge e^b\wedge e^c\wedge e^d
 +S_{\rm m}[e,\omega,\Theta].
 \label{eq:palatini-action}
\end{equation}
Define the energy-momentum three-form by
$\delta_eS_{\rm m}=\int\Sigma_a\wedge\delta e^a$ and spin current as in the
companion torsion note.

\begin{theorem}[First-order field equations]
Independent variation of Eq.~\eqref{eq:palatini-action} gives
\begin{align}
 \epsilon_{abcd}\,e^b\wedge R^{cd}
 -\frac{\Lambda}{3}\epsilon_{abcd}\,e^b\wedge e^c\wedge e^d
 &= -2\kappa\,\Sigma_a,
 \label{eq:tetrad-eom}\\
 \epsilon_{abcd}\,T^a\wedge e^b&=\kappa\tau_{cd}.
 \label{eq:connection-eom}
\end{align}
When $\tau_{ab}=0$, the second equation gives the Levi--Civita connection and
the first is equivalent to
\begin{equation}
 \boxed{G_{\mu\nu}+\Lambda g_{\mu\nu}=\kappa T_{\mu\nu}.}
 \label{eq:einstein}
\end{equation}
\end{theorem}

\begin{remark}
This is standard tetrad/Palatini gravity.  It proves what follows from the
minimal action; it does not prove that Eq.~\eqref{eq:palatini-action} has
already been derived from the deeper UBT invariant.
\end{remark}

\section{Four-dimensional uniqueness}
\begin{theorem}[Lovelock consequence in four dimensions]
Assume the classical metric equation is a local natural symmetric rank-two
tensor $\mathcal E_{\mu\nu}[g]$ which
\begin{enumerate}
 \item depends on the metric and at most its first two derivatives;
 \item is divergence free identically;
 \item contains no additional propagating geometric fields in the branch;
 \item is of the standard second-order/Lovelock class.
\end{enumerate}
Then in four dimensions
\begin{equation}
 \boxed{\mathcal E_{\mu\nu}=aG_{\mu\nu}+bg_{\mu\nu}.}
 \label{eq:lovelock-form}
\end{equation}
After fixing normalization and writing $\Lambda=b/a$, the sourced equation is
Einstein's equation with cosmological constant.
\end{theorem}

\begin{corollary}[UBT low-energy selection]
If the on-shell UBT tetrad branch is local, diffeomorphism and local-Lorentz
invariant, torsion free at macroscopic scales, second order, and contains no
additional light geometric fields, then its metric dynamics cannot be an
arbitrary alternative tensor equation: it is Einstein--$\Lambda$ up to the
overall coupling.
\end{corollary}

\section{What this does not derive}
The theorem does not determine from UBT:
\begin{itemize}
 \item why the effective action is local and in the Lovelock class;
 \item why higher-curvature or nonlocal terms are absent or suppressed;
 \item the values and signs of $\kappa$ and $\Lambda$;
 \item the exact matter stress tensor and paired spin current;
 \item the on-shell solution of $E_\mu=D_\mu\Theta/\sqrt{\mathcal N_0}$.
\end{itemize}

\section{Status}
\begin{tcolorbox}[colback=green!6!white,colframe=green!55!black,title=Defensible status]
\textbf{GAP-10D-PALATINI: CLOSED CONDITIONALLY [L1].}
The minimal first-order tetrad action yields the Cartan equation and
Einstein--$\Lambda$ dynamics.

\medskip
\textbf{GAP-10D-UNIQUENESS: CLOSED CONDITIONALLY [L1].}
Under the four-dimensional Lovelock/low-energy assumptions, the only metric
equation is $aG_{\mu\nu}+bg_{\mu\nu}=\kappa T_{\mu\nu}$.

\medskip
\textbf{GAP-10D: NARROWED, NOT FULLY CLOSED.}
The remaining UBT task is to derive the minimal low-energy action or the
Lovelock hypotheses, its coefficients and matter terms from the canonical
single-field theory.  The conditional infrared endpoint is no longer
ambiguous.
\end{tcolorbox}

\begin{thebibliography}{9}
\bibitem{lovelock1971} D.~Lovelock, ``The Einstein Tensor and Its
Generalizations,'' \emph{J. Math. Phys.} \textbf{12}, 498--501 (1971).
\bibitem{hehl1976} F.~W.~Hehl, P.~von der Heyde, G.~D.~Kerlick, and J.~M.~Nester,
``General Relativity with Spin and Torsion: Foundations and Prospects,''
\emph{Rev. Mod. Phys.} \textbf{48}, 393--416 (1976).
\end{thebibliography}
\end{document}
